\subsection{Banachscher Fixpunktsatz}

\begin{definition}{}
Sei \( (M, d) \) ein metrischer Raum. Eine Abbildung \( \varphi: M \to M \) heißt \emph{kontrahierend}, falls es ein \( \lambda \in [0,1) \) gibt mit


\[
d(\varphi(x), \varphi(y)) \leq \lambda d(x,y) \quad \forall x,y \in M.
\]


\end{definition}

\begin{bemerkungbox}
Eine kontrahierende Abbildung ist immer Lipschitz-stetig mit Lipschitz-Konstante \( \lambda < 1 \).
\end{bemerkungbox}

\begin{satz}{}
Sei \( (M, d) \) ein vollständiger metrischer Raum und sei \( \varphi: M \to M \) kontrahierend. Dann besitzt \( \varphi \) genau einen Fixpunkt \( z \), d.h.


\[
\varphi(z) = z.
\]


Außerdem konvergiert für jedes \( x_0 \in M \) die Folge


\[
x_{n+1} = \varphi(x_n)
\]


gegen \( z \).
\end{satz}

\begin{beweisbox}
Die Existenz des Fixpunkts folgt aus der Konstruktion einer Cauchy-Folge \( (x_n) \), die durch wiederholtes Anwenden von \( \varphi \) definiert ist. Da \( M \) vollständig ist, konvergiert diese Folge zu einem Grenzwert \( z \), der sich als Fixpunkt zeigt:


\[
\varphi(z) = \lim_{n \to \infty} \varphi(x_n) = z.
\]


Für die Eindeutigkeit sei \( z, w \) Fixpunkte mit \( \varphi(z) = z \) und \( \varphi(w) = w \). Dann folgt


\[
d(z, w) = d(\varphi(z), \varphi(w)) \leq \lambda d(z,w),
\]


was nur für \( d(z,w) = 0 \) möglich ist, also \( z = w \).
\end{beweisbox}

\begin{beispielbox}
\begin{enumerate}
    \item Sei \( \varphi(x) = \frac{x}{2} \) auf \( \mathbb{R} \). Diese Funktion ist kontrahierend mit \( \lambda = \frac{1}{2} \) und besitzt genau einen Fixpunkt \( z = 0 \).
    \item Newtons Verfahren zur Nullstellenbestimmung kann als Anwendung des Fixpunktsatzes interpretiert werden. Die Fixpunkte einer Iterationsfunktion sind die Lösungen der Gleichung \( f(x) = 0 \).
\end{enumerate}
\end{beispielbox}
